\section{问题重述与研究路线}
题目要求在算力约束下提高大语言模型能力。问题一使用附件 A 的文本质量信号、配比及逐目标 Loss，评价数据质量、处理指标冲突并建立配方决策；问题二使用附件 B 的训练日志与质量实验建立规模、数据量、质量和配比的统一预测关系；问题三在预算及上下文成本下联合配置参数量、token、质量与配比；问题四结合附件 C 的评测、资源、上下文及 Loss--Benchmark 信息，讨论能力演进、技术因素与算力增长放缓的情景。

四问并非同一批模型的完整控制实验。本文先按观测、半合成、插值、估算和参考资料固定数据角色，再建立各附件自身可识别的关系。跨附件仅通过显式条件接口连接；没有成对观测的 $Q$ 或 Loss 坐标不补造标定参数。每问同时报告模型内可计算结果、可用验证和超出数据证据的边界。主要数据字段、单位、模型代码及实验清单分别见各问结果目录。
